\section{Verification and Quality Assurance}\label{sec:implementation-verification}

To ensure the correctness of the compilation pipeline and the validity of the generated MIDI artifacts, a
comprehensive testing suite was implemented using the \textit{Google Test} framework.
The testing strategy follows a tiered approach, covering individual units, module boundaries, and end-to-end
integration.

\subsection{Unit Testing}

Unit tests focus on the smallest verifiable components of the compiler.
This is particularly important for the musical domain logic, where small errors in calculation can lead to
significant musical dissonance.
Key test cases include:
\begin{itemize}
  \item \textbf{Pitch Math:} Verifying that pitch literals (e.g., \texttt{A4}) map correctly to MIDI note
    numbers (e.g., 69) and that octave offsets are handled according to the standard.
  \item \textbf{Diagnostic Reporting:} Ensuring that the \texttt{DiagnosticsEngine} correctly captures and
    formats errors across different stages of the pipeline.
  \item \textbf{Lexer Accuracy:} Stress-testing the Flex regex patterns with edge cases, such as multiple
    consecutive accidentals or varying whitespace.
\end{itemize}

\subsection{Integration and Regression Testing}

Integration tests verify that the different stages of the compiler work together seamlessly.
These tests utilize a "Golden File" methodology:
\begin{enumerate}
  \item A set of representative DSL source files (located in \texttt{tests/data/}) is compiled.
  \item The resulting IR is checked for temporal consistency (e.g., ensuring no notes overlap unless
    explicitly requested).
  \item The final MIDI binary is compared against a previously validated "golden" version to detect
    unexpected changes in behavior.
\end{enumerate}

Regression tests were instrumental during the refactoring process (e.g., when decomposing the monolithic lowerer).
By maintaining a suite of complex compositions—such as the \texttt{fur\_elise.dsl} example—the compiler's
output could be verified for bit-perfect consistency after every architectural change.

\subsection{Integration in the Development Lifecycle}

The testing suite is integrated into the build system via CMake's \texttt{CTest} module.
This allows for one-command verification of the entire system (\texttt{make test}).
By enforcing a high level of test coverage, particularly in the lowering and backend stages, the
implementation provides high confidence that the generated compositions are mathematically and musically accurate.
This rigorous verification process is a cornerstone of the thesis, transforming a creative tool into a
reliable software system.
